\subsection{An exterior spectral estimate}

The outward cubic, quadratic and degree-loss passages share a scalar
graph equation. The quadratic coefficient determines a power in its
first variation, while the remaining terms control the error in its
logarithmic derivatives. We use logarithmic height as the independent
variable to estimate these derivatives directly, including the moving
endpoints.

\begin{lemma}[Exterior spectral estimate]\label{deg:lem-spectral-exterior}
Fix a terminal cut \(d>0\), a small lip constant \(\ell_0>0\), compact
bounds on \(A,C\), and positive constants \(c_0,C_0\). Put
\[
 \begin{gathered}
 M=tx^3+\mu_2x^2+\mu_1x+\mu_0,\qquad B=1+b+tx,\\
 Q=1-A(v+tx)+C(v+tx)^2,\quad \zeta=tC(v+tx),\quad |t|\le1.
 \end{gathered}
\]
All coefficients are fixed under \(E=a\partial_a\). On an enlarged
open interval of \(0<a<d/2\), suppose, pointwise in \(a\), that
some \(0<\rho\le\rho_0\) satisfies
\begin{equation}\label{deg:spectral-exterior-hypotheses}
 \begin{gathered}
 c_0\rho\le m(x):=M(x)/x^2\le C_0\rho\quad(a\le x\le2a),\\
 |b|+|v|+|\mu_2|+|t|a+|\mu_1|/a+|\mu_0|/a^2\le C_0\rho .
 \end{gathered}
\end{equation}
For sufficiently small \(|t|d,\rho_0\), depending only on these
fixed bounds, the graph
\begin{equation}\label{deg:spectral-exterior-graph}
 z_x=\frac{z(xB+\zeta z)}{M+Qz},\qquad z(a)=\ell_0M(a)
\end{equation}
reaches \(x=d\) with positive height and denominator. An ordinary
\(C^3\) observer \(\chi\), with \(\chi'>0\), is appended on the
open preimage of its prescribed buffered terminal image. Set
\[
 \nu=\frac{2\mu_2}{1+b},\quad\lambda=2\nu,\quad e=2-\lambda>3/2,
 \quad\Theta=|t|a+|\mu_1|/a+|\mu_0|/a^2,\quad L_\rho=1+|\log\rho|,
\]
\[
 H_\chi=\sup\{|z(\log\chi')'|
       +|z(\log\chi')'+z^2(\log\chi')''|\},\qquad
 \mathcal T=|t|d+|\mu_2|+|\mu_1|/d+|\mu_0|/d^2+H_\chi ,
\]
where the supremum is on that buffer and is finite.
For any coefficient-only \(\sigma>0\), use
\(\beta=-a^2/(2\sigma^2)\) and \(\Gamma(\beta)=\chi(z_d)\).
Then \(\Gamma'>0\), and
\begin{equation}\label{deg:spectral-exterior-rows}
 |E\log\Gamma'+\lambda|+|E^2\log\Gamma'|
 \le C\{\Theta L_\rho+\mathcal T(a/d)^e\}.
\end{equation}
The same rows hold after a positive coefficient-only affine calibration.
For intrinsic terminal height take \(\chi=\mathrm{id}\), so \(H_\chi=0\).
No reciprocal bound on \(\sigma\) or bound on \(\lambda\log a\) is used.
\end{lemma}
\begin{proof}
Set \(W=z/x^2\), \(\tau=\log(x/a)\), \(N=B+\zeta xW\), and
\(D=m+QW\). Then \(W_\tau=F=W(N/D-2)\). On \([a,2a]\),
\(N,Q\) are positive units and \(D\asymp\rho+W\) for bounded \(W\).
Below a small fixed \(w_0\), \(F\ge cW/(\rho+W)\). Starting at
\(W_l=\ell_0m(a)\), the time to reach \(w_0\) is at most
\[
 C\rho\log\frac{w_0}{\ell_0c_0\rho}+Cw_0<\log2,
\]
by choosing \(w_0\) first and then \(\rho_0\). At the boundaries
\(w_0<1/2\), \(w_1=3/4\), the formula \(F(x,w)=1-2w+O_w(\rho)\)
makes the strip invariant up to \(2a\). Beyond it,
\(|m(x)|\le|t|d+C\rho\ll w_0^2\) preserves the strip, gives
positive-unit \(D\) and \(F_W\le-c<0\). Thus the graph reaches
\(d\), with \(c\rho\le W\le w_1\), even if \(m\) changes sign.

The bounded strong graph with \(b=v=\mu_j=0\), keeping \(t,A,C\),
is \(W_*=1/2+O(|t|x)\), by the integral from zero for
\(xU'=-2U+O(|t|x)\), \(U=W-1/2\). Strip contraction gives
\begin{equation}\label{deg:spectral-exterior-window}
 |W(d)-W_*(d)|\le C\rho+C(a/d)^c.
\end{equation}
Where this comparison applies, the original buffer margin and the
permitted cutoff choices put the terminal height in the observer's
domain. Elsewhere we retain the original terminal admission condition.

Put \(r=\log x\), \(\eta=\log z\). At fixed height, \(W_r=-2W\)
and \(r_\eta=V=D/N\). The raw cocycle
\eqref{deg:scalar-raw-cocycle} has the exact density split
\begin{equation}\label{deg:spectral-exterior-density}
 \begin{gathered}
 \frac{2m+xm'+(xQ'+\zeta x)W}{N}=\nu+J,\\
 J=\frac{(3-\nu)tx+\mu_1/x+
          tx[-A+(3-\nu)C(v+tx)]W}{N}.
 \end{gathered}
\end{equation}
Here \(2m+xm'=3tx+2\mu_2+\mu_1/x\); \(\mu_0\) cancels only
in the density. Let
\(I=\int_{\eta_l}^{\eta_T}J\,d\eta\), with
\(\eta_l=\log(\ell_0M(a))\). The near height length is \(O(L_\rho)\);
on the strip \(d\eta\le C\,d\log x\). Thus every partial integral
is bounded by \(C\{\Theta L_\rho+|t|d+|\mu_1|/a\}\).

Use \eqref{deg:scalar-inverse-endpoints} with
\(P_j=E^jr|_\eta\), \(l_j=E^j\eta_l\), \(t_j=E^j\eta_T\),
\(Er_l=1\). The polynomial gives bounded \(l_1,l_2\),
\(V_l,V_{r,l}=O(\rho)\), and \(P_{1,l}=1+O(\rho)>0\).
For \(\mathscr B=EW_l-F_l\), the exact first cocycle gives
\begin{equation}\label{deg:spectral-exterior-inverse-decay}
 \begin{gathered}
 P_1=\frac{-\mathscr B(1+\ell_0Q(a))}{\ell_0N}
       (a/x)^e\left(\frac W{\ell_0m(a)}\right)^\nu e^I,\\
 0<P_1\le C(a/x)^e,\qquad |P_2|\le CP_1 .
 \end{gathered}
\end{equation}
The initial factor equals \(B(a)K(a)>0\), with \(K\) from
\eqref{deg:scalar-polynomial-lip-slope}. Since \(|\nu|\le C\rho\)
and \(c\rho\le W\le w_1\), its height ratio to power \(\nu\) is a
unit for either sign, proving the first bound without bounding
\(a^{-\lambda}\). Near the lip, fixed-height differentiation gives
\(|V_{rr}|\le C(\rho+W)\), whereas \(d\eta=N/D\,d\log x\).
On the strip all these factors are bounded. Hence
\[
 |V_{rr}|\,d\eta\le C\,d\log x,\qquad
 (P_2/P_1)_\eta=V_{rr}P_1.
\]
Its integral is bounded by \(C\int_a^d(a/x)^e\,d\log x\), proving
the second bound. The two fixed-terminal rows then give
\(|t_1|+|t_2|\le C(a/d)^e\).

Rational differentiation at fixed height gives
\(|J_r|+|J_{rr}|\le C(|t|x+|\mu_1|/x)\). In the two full Leibniz
rows \eqref{deg:scalar-density-endpoints}, the lower terms cost
\(C\Theta\), the near integrals cost \(C\Theta L_\rho\), and
\[
 \int_{2a}^d(|t|x+|\mu_1|/x)(a/x)^e\,d\log x
 \le C(|t|a+|\mu_1|/a).
\]
Every upper term retains \((a/d)^e\). Therefore
\begin{equation}\label{deg:spectral-exterior-integral-rows}
 |EI|+|E^2I|
 \le C\{\Theta L_\rho+(|t|d+|\mu_1|/d)(a/d)^e\}.
\end{equation}
The exact polynomial-lip formula becomes
\begin{equation}\label{deg:spectral-exterior-slope}
 \Gamma'=\sigma^2a^{-\lambda}
 B(a)K(a)\ell_0^{-\nu}m(a)^{-\nu}
 \frac{W_T}{m(d)+Q(d)W_T}z_T^\nu\chi'(z_T)e^I .
\end{equation}
Both \(E\) rows of \(\log(BK)\) are \(O(\Theta)\). Also
\(E m,E^2m=O(\Theta)\), \(m(a)\asymp\rho\), and
\(|\nu|,\Theta\le C\rho\), so those of \(-\nu\log m(a)\) cost
\(C\Theta\). At fixed \(d\), the first two height-logarithmic rows
of \(\log[W_T/(m(d)+Q(d)W_T)]\) cost \(C|m(d)|\), and those of
\(\nu\log z_T\) cost \(C|\nu|\). The observer rows are precisely
those defining \(H_\chi\). Their terminal chain rules and the integral
bounds prove \eqref{deg:spectral-exterior-rows}, since
\(a^{-\lambda}\) contributes \(-\lambda\) and zero to the first and
second Euler rows.
\end{proof}
